% SEGMENT OLP-0606-S01
% Part: methods
% Chapter: proofs
% Section: inference-patterns

\documentclass[../../../include/open-logic-section]{subfiles}

\begin{document}

\olfileid{mth}{prf}{pat}

\olsection{உய்த்தறிதல் வடிவங்கள்}

நிறுவல்கள் தனித்தனி உய்த்தறிதல்களால் அமைகின்றன. ஓர்
உய்த்தறிதலைச் செய்யும்போது, ``ஆக'', ``எனவே'', ``ஆதலால்''
போன்ற சொல்லால் அதை வழக்கமாகக் குறிக்கிறோம். அந்த உய்த்தறிதல்,
நிறுவலில் ஏற்கெனவே கிடைக்கும் ஓரிரு உண்மைகளை அடிக்கடி
சார்ந்திருக்கும். அவை நாம் எடுத்துக்கொண்ட எடுகோள்களாகவோ,
முன்னர் உய்த்தறிந்த முடிவுகளாகவோ இருக்கலாம். தெளிவுக்காக
அவற்றைக் குறியிட்டு, புதிய உய்த்தறிதலில் எந்த முந்தைய
கூற்றுகளைப் பயன்படுத்துகிறோம் என்று சொல்லலாம். முந்தையவற்றிலிருந்து
புதிய முடிவு \emph{ஏன்} பின்தொடர்கிறது என்பதற்கான விளக்கமும்
உய்த்தறிதலில் அடிக்கடி இருக்கும். நிறுவல்களில் சில பொதுவான
உய்த்தறிதல் வடிவங்கள் மீண்டும் மீண்டும் பயன்படுகின்றன; சிலவற்றை
இங்கே பார்ப்போம். தொகுத்தறிதல் வழி நிறுவல் போன்ற பிற வடிவங்கள்
சற்றுச் சிக்கலானவை; அவை பின்னர் விவாதிக்கப்படும்.

% SEGMENT OLP-0606-S02
ஒரு வடிவத்தை நாம் ஏற்கெனவே பார்த்தோம்: வரையறையை
விரித்துரைப்பது அல்லது பயன்படுத்துவது. அதை விரிக்கும்போது,
வரையறுக்கப்படும் சொல்லைக் கொண்ட கூற்றை, அதை வரையறுக்கும்
பகுதியால் மீண்டும் கூறுகிறோம். எடுத்துக்காட்டாக, நிறுவலின்
போது $D = E$ என்பதை (a) என்று ஏற்கெனவே நிறுவியதாகக்
கொள்க. அப்போது கணங்களுக்கான $=$ இன் வரையறையைப்
பயன்படுத்தி, ``ஆக, (a) இலிருந்தும் வரையறையின்படியும்,
$D$ இன் ஒவ்வோர் !!{element} உம் $E$ இன்
!!a{element} ஆகவும், மறுதிசையிலும் இருக்கும்'' என்று
உய்த்தறியலாம்.

% SEGMENT OLP-0606-S03
சற்றுக் குழப்பமூட்டும் விதமாக, உய்த்தறிதலுக்கான நியாயத்தை
அதைச் செய்யும் இடத்தில் எழுதாமல், அதற்கு முன்பே எழுதுவது
வழக்கம். $D = E$ என்பதை இன்னும் நிறுவவில்லை; அதையே
நிறுவ விரும்புகிறோம் என்க. $D = E$ என்பதே நோக்கிய முடிவானால்,
வரையறையைப் பயன்படுத்தி இலக்கை மீண்டும் கூறலாம்:
$D = E$ என்பதை நிறுவ, $D$ இன் ஒவ்வோர் !!{element}
உம் $E$ இன் !!a{element} ஆகவும், மறுதிசையிலும்
இருக்கிறது என்று நிறுவ வேண்டும். ஆகவே நிறுவலின் வடிவம்
இவ்வாறு அமையும்: (a) $D$ இன் ஒவ்வோர் !!{element}
உம் $E$ இன் !!a{element} என்று நிறுவுக; (b) $E$ இன்
ஒவ்வோர் !!{element} உம் $D$ இன் !!a{element}
என்று நிறுவுக; (c) ஆக (a), (b) ஆகியவற்றிலிருந்தும்
$=$ இன் வரையறையின்படியும் $D = E$. ஆனால் வழக்கமாக
இவ்வாறு எழுதுவதில்லை. அதற்குப் பதிலாகப் பின்வருமாறு
எழுதலாம்:
\begin{quote}
$D = E$ என்று காட்ட வேண்டும். $=$ இன் வரையறையின்படி,
இது $D$ இன் ஒவ்வோர் !!{element} உம் $E$ இன்
!!a{element} ஆகவும், மறுதிசையிலும் இருக்கிறது என்று
காட்டுவதற்குச் சமம்.

(a) \dots ($D$ இன் ஒவ்வோர் !!{element} உம் $E$ இன்
!!a{element} என்று காட்டும் நிறுவல்) \dots

(b) \dots ($E$ இன் ஒவ்வோர் !!{element} உம் $D$ இன்
!!a{element} என்று காட்டும் நிறுவல்) \dots
\end{quote}

% SEGMENT OLP-0606-S04
\subsection{இணைப்பைப் பயன்படுத்துதல்}

மிக எளிய உய்த்தறிதல் வடிவங்களில் ஒன்று, ஓர் இணைப்புக்
கூற்றின் இரு கூறுகளில் ஒன்றை முடிவாகப் பெறுவது. வேறுவிதமாகச்
சொன்னால், $p$ உம் $q$ உம் மெய் என்று எடுத்துக்கொண்டாலோ
ஏற்கெனவே நிறுவியிருந்தாலோ, $p$ ஐ (அதேபோல் $q$ ஐயும்)
உய்த்தறியலாம். இந்த உய்த்தறிதல் மிகவும் அடிப்படையானதால்,
அதை வெளிப்படையாகக் குறிப்பிடாமல் இருப்பது வழக்கம்.
எடுத்துக்காட்டாக, $D = E$ இன் வரையறையை விரித்தால்,
$D$ இன் ஒவ்வோர் !!{element} உம் $E$ இன்
!!a{element} ஆகவும், மறுதிசையிலும் இருப்பதைப் பெறுகிறோம்.
இதிலிருந்து $E$ இன் ஒவ்வோர் !!{element} உம் $D$ இன்
!!a{element} என்று முடிக்கலாம்; இது ``மறுதிசையிலும்''
என்பதன் பகுதி.

% SEGMENT OLP-0606-S05
\subsection{இணைப்புக் கூற்றை நிறுவுதல்}

சில நேரம் நிறுவ வேண்டியது இணைப்புக் கூற்றாக இருக்கும்:
``$p$ உம் $q$ உம் மெய் என்று நிறுவுக''. அப்போது
இரு வேலைகளைச் செய்ய வேண்டும்: முதலில் $p$ ஐ நிறுவி,
பிறகு $q$ ஐ நிறுவ வேண்டும். தெளிவுக்காக நிறுவலை
இரு பகுதிகளாகப் பிரித்துக் குறியிடலாம். தொடக்கக் குறிப்புகளில்
தாளின் மேற்பகுதியில் ``(1) $p$ ஐ நிறுவுக'' என்றும்
நடுப்பகுதியில் ``(2) $q$ ஐ நிறுவுக'' என்றும் எழுதலாம்.
இணைப்புக் கூற்றை நிறுவச் சொல்லாமலேயே, உங்கள் நிறுவலுக்கு
அதைக் காட்ட வேண்டியிருக்கலாம். எடுத்துக்காட்டாக, $D = E$
என்பதை நிறுவும்போது $=$ இன் வரையறையை விரித்த பின்,
$D$ இன் ஒவ்வோர் !!{element} உம் $E$ இன்
!!a{element} ஆகும் \emph{மேலும்} $E$ இன் ஒவ்வோர்
!!{element} உம் $D$ இன் !!a{element} ஆகும்
என்று நிறுவ வேண்டியிருக்கும்.

% SEGMENT OLP-0606-S06
\subsection{பிரிப்புக் கூற்றை நிறுவுதல்}

நிறுவ வேண்டியது ஒரு பிரிப்புக் கூற்று, அதாவது ``$p$
அல்லது $q$'' என்ற வடிவிலுள்ள கூற்று எனில், அதன்
இரு கூறுகளில் ஒன்று மெய் என்று காட்டுவது போதும்.
ஆனால் உங்கள் தேற்றத்தின் எடுகோள்களிலிருந்து ஒரு கூறு
நேரடியாகப் பின்தொடர்வது நடைமுறையில் கிட்டத்தட்ட நிகழாது.
பெரும்பாலும் எடுகோள்களே பிரிப்புக் கூற்றாக இருக்கும்;
அல்லது ஒரு வகையைச் சேர்ந்த எல்லாப் பொருள்களுக்கும்
இரு பண்புகளில் ஒன்று இருப்பதைக் காட்டுவீர்கள்: சிலவற்றுக்கு
ஒரு பண்பும் மற்றவற்றுக்கு மற்றொரு பண்பும் இருக்கும்.
இத்தகைய இடங்களில் நிகழ்வுகளைப் பிரித்து நிறுவுவது
பயனுள்ளது (கீழே பார்க்கவும்).

% SEGMENT OLP-0606-S07
\subsection{நிபந்தனைக் கூற்றை நிறுவுதல்}

பல தேற்றங்கள் நிபந்தனை வடிவில் இருக்கும்: $p$ மெய்யானால்
$q$ உம் மெய் என்று காட்ட வேண்டும். இவற்றுக்கான
நிறுவலை அமைப்பது எளிது: நிபந்தனையின் முற்பகுதியை,
இங்கு $p$ ஐ, எடுகோளாக எடுத்துக்கொண்டு அதிலிருந்து
முடிவான $q$ ஐ நிறுவுக. ஆகவே ``$p$ எனில் $q$''
என்று தேற்றம் சொன்னால், ``$p$ என்க'' என்று நிறுவலைத்
தொடங்கி, இறுதியில் $q$ ஐ நிறுவியிருக்க வேண்டும்.

நிபந்தனைக் கூற்றை வெவ்வேறு வழிகளில் கூறலாம். ``$p$
எனில் $q$'' என்பதற்குப் பதிலாக, ``$p$ என்பது $q$
எனில் மட்டுமே'', ``$p$ எனில் $q$'', அல்லது ``$p$
என்ற நிபந்தனையில் $q$'' என்றும் தேற்றம் கூறலாம்.
இவையனைத்தும் ஒரே பொருள் தருகின்றன: $p$ ஐ எடுகோளாக
எடுத்து அதிலிருந்து $q$ ஐ நிறுவ வேண்டும். ``$p$
எனில், எனில் மட்டுமே, $q$'' என்ற இருதிசை நிபந்தனை
உண்மையில் இரண்டு நிபந்தனைகளின் இணைப்பு: $p$ எனில்
$q$; $q$ எனில் $p$. எனவே இருமுறை நிபந்தனை
நிறுவல் போதும்; ஒவ்வொரு திசைக்கும் ஒன்று. சில நேரம்
பல ``எனில், எனில் மட்டுமே'' கூற்றுகளைச் சங்கிலியாக
இணைத்து $p$ இல் தொடங்கி $q$ இல் முடிப்பதாலும்
இருதிசை நிபந்தனையை நிறுவலாம். அப்போது ஒவ்வொரு
படியும் உண்மையிலேயே இருதிசை நிபந்தனை என்று உறுதிசெய்ய
வேண்டும்.

% SEGMENT OLP-0606-S08
\subsection{அனைத்தையும் பற்றிய கூற்றுகள்}

அனைத்தையும் பற்றிய கூற்றைப் பயன்படுத்துவது எளிது:
ஒவ்வொன்றுக்கும் மெய்யானது, குறிப்பிட்ட ஒவ்வொன்றுக்கும்
மெய்யாகும். எடுத்துக்காட்டாக, நிறுவலின் எடுகோள்
$A \subseteq B$ எனில், $\subseteq$ இன் வரையறையை
விரித்தால் ஒவ்வொரு $x \in A$ க்கும் $x \in B$ என்று
பொருள். எனவே $z \in A$ என்று ஏற்கெனவே தெரிந்தால்,
$z \in B$ என்று முடிக்கலாம்.

அனைத்தையும் பற்றிய கூற்றை நிறுவுவது சற்றுக் கடினமாகத்
தோன்றலாம். வழக்கமாக அவை ``$x$ க்கு $P$ என்ற
பண்பு இருந்தால் $Q$ என்ற பண்பும் இருக்கும்'' அல்லது
``$P$ பண்புள்ள எல்லாவற்றுக்கும் $Q$ பண்பு உண்டு'' என்ற வடிவில் இருக்கும்.
எப்போதும் இதே வடிவில் சரியாக அமையாமல் இருக்கலாம்;
எதை நிறுவக் கேட்கிறார்கள் என்று கண்டறியப் பயிற்சி தேவை.
ஆனால் ஒரு பண்புள்ள எல்லாப் பொருள்களுக்கும் மற்றொரு
பண்பு உள்ளது என்பதை நிறுவ வேண்டியிருக்கிறது.

% SEGMENT OLP-0606-S09
அப்படிப்பட்ட கூற்றை நிறுவ, முதல் பண்பைக் கொண்ட
பொருள்களுக்குப் பெயர்களையோ மாறிகளையோ அறிமுகப்படுத்தி,
அவற்றுக்கு மற்றப் பண்பும் உள்ளது என்று காட்ட வேண்டும்.
\emph{எல்லா} $P$ பண்புடையவற்றைப் பற்றியும் ஒன்றை நிறுவ,
\emph{தன்னிச்சையான} ஒரு $P$ பற்றியே அதை நிறுவ வேண்டும்
என்று சொல்லலாம். அறிமுகப்படுத்தப்படும் பெயர் அத்தகைய
தன்னிச்சையான $P$ க்கான பெயராகும். வழக்கமாகத் தனி
எழுத்துகளையே பயன்படுத்துகிறோம்; குறிப்பிட்ட மரபுகளும்
உள்ளன: இயல் எண்களுக்கு $n$, !!{formula}s க்கு $!A$,
கணங்களுக்கு $A$, சார்புகளுக்கு $f$ போன்றவை.

நிறுவல் முழுவதும் பொதுத்தன்மையைப் பாதுகாப்பதே நுணுக்கம்.
ஒரு தன்னிச்சையான பொருளுக்கு (``$x$'') $P$ பண்பு
உள்ளது எனத் தொடங்கி, வரையறைகளையும் அனுமதிக்கப்பட்ட
எடுகோள்களையும் மட்டுமே கொண்டு $x$ க்கு $Q$ பண்பும்
உள்ளது என்று காட்டுக. $P$ உள்ளது என்பதைத் தவிர $x$
குறிப்பாக என்னவென்று நிர்ணயிக்கவில்லை. ஆகவே $P$
உள்ள எல்லாவற்றுக்கும் $Q$ உள்ளது என்று கூறலாம்.
சுருக்கமாக, $P$ பண்புள்ள \emph{எல்லா} பொருள்களுக்குமான
இடநிறைவாக $x$ இருக்கிறது.

\begin{prop}
  எல்லாக் கணங்கள் $A$, $B$ க்கும் $A \subseteq A \cup B$.
\end{prop}

\begin{proof}
  $A$, $B$ தன்னிச்சையான கணங்களாக இருக்கட்டும்.
  $A \subseteq A \cup B$ என்பதைக் காட்ட வேண்டும்.
  $\subseteq$ இன் வரையறையின்படி, ஒவ்வொரு $x$ க்கும்
  $x \in A$ எனில் $x \in A \cup B$ என்று காட்ட வேண்டும்.
  எனவே $x \in A$ என்பதை, $A$ இன் தன்னிச்சையான
  ஓர் !!{element} ஆக எடுத்துக்கொள்க. $x \in A \cup B$
  என்று காட்ட வேண்டும். $x \in A$ என்பதால்,
  $x \in A$ அல்லது $x \in B$. ஆகவே
  $x \in \Setabs{x}{x \in A \lor x \in B}$. இது
  $\cup$ இன் வரையறையின்படி $x \in A \cup B$
  என்பதையே குறிக்கிறது.
\end{proof}

% SEGMENT OLP-0606-S10
\subsection{நிகழ்வுகளைப் பிரித்து நிறுவுதல்}

ஒரு பிரிப்புக் கூற்றை எடுகோளாகக் கொண்டிருக்கிறீர்கள்
அல்லது ஏற்கெனவே முடித்திருக்கிறீர்கள் என்க: $p$ அல்லது
$q$ மெய். நீங்கள் $r$ ஐ நிறுவ விரும்புகிறீர்கள்.
இதை இரு படிகளில் செய்க: முதலில் $p$ மெய் என்று
எடுத்துக்கொண்டு $r$ ஐ நிறுவுக; பிறகு $q$ மெய் என்று
எடுத்துக்கொண்டு மீண்டும் $r$ ஐ நிறுவுக. இரு மாற்றுகளில்
ஒன்று மெய் என்று எடுகோளாகக் கொண்டிருக்கிறோம் அல்லது
அறிந்திருக்கிறோம் என்பதால் இது செயல்படுகிறது. எந்த
மாற்றாக இருந்தாலும் $r$ மெய் என்பதற்கு இந்த இரு
படிகளும் போதுமானவை. இரண்டும் மெய்யானால், $r$ மெய்
என்பதற்கு ஒன்று அல்ல, இரண்டு காரணங்கள் கிடைக்கும்.
$p$, $q$ இரண்டும் மெய் என்ற தனி வழக்கிலும் $r$
மெய் என்று மீண்டும் நிறுவத் தேவையில்லை. நாம் செய்வதைக்
குறிக்க ``நிகழ்வுகளைப் பிரிக்கிறோம்'' என்று அறிவிக்கலாம்.
எடுத்துக்காட்டாக, $x \in B \cup C$ என்று தெரியும் என்க.
$B \cup C$ என்பது $\Setabs{x}{x \in B \text{ அல்லது } x \in C}$
என்று வரையறுக்கப்படுகிறது. எனவே வரையறையின்படி
$x \in B$ அல்லது $x \in C$. இதிலிருந்து $x \in A$
என்று நிறுவ, முதலில் $x \in B$ என எடுத்துக்கொண்டு
அதிலிருந்து $x \in A$ என்று காட்டுவோம்; பிறகு
$x \in C$ என எடுத்துக்கொண்டு மீண்டும் $x \in A$
என்று காட்டுவோம். எடுகோளுக்குக் கீழே ``நிகழ்வுகளைப்
பிரிக்கிறோம்'' என்றும், அதன் கீழே ``நிகழ்வு (1):
$x \in B$'' என்றும், தாளின் நடுவில் ``நிகழ்வு (2):
$x \in C$'' என்றும் எழுதலாம். பின்னர் தாளின்
மேற்பாதியிலும் கீழ்ப்பாதியிலும் உரிய நிறுவல்களை நிரப்பலாம்.

% SEGMENT OLP-0606-S11
நிறுவ வேண்டியதே பிரிப்புக் கூற்றாக இருந்தால், நிகழ்வுகளைப்
பிரித்தல் குறிப்பாகப் பயனுள்ளது. இதோ எளிய எடுத்துக்காட்டு:

\begin{prop}
$B \subseteq D$ மற்றும் $C \subseteq E$ என்க. அப்போது
$B \cup C \subseteq D \cup E$.
\end{prop}

\begin{proof}
  (a) $B \subseteq D$ என்றும் (b) $C \subseteq E$
  என்றும் எடுத்துக்கொள்க. வரையறையின்படி, எந்த $x \in B$
  உம் $\in D$ (c); எந்த $x \in C$ உம் $\in E$ (d).
  $B \cup C \subseteq D \cup E$ என்று காட்ட,
  $x \in B \cup C$ எனில் $x \in D \cup E$ என்று
  காட்ட வேண்டும் ($\subseteq$ இன் வரையறையின்படி).
  $x \in B \cup C$ எனில், எனில் மட்டுமே,
  $x \in B$ அல்லது $x \in C$ ($\cup$ இன்
  வரையறையின்படி). அதேபோல் $x \in D \cup E$
  எனில், எனில் மட்டுமே, $x \in D$ அல்லது $x \in E$.
  ஆகவே ஒவ்வொரு $x$ க்கும், $x \in B$ அல்லது
  $x \in C$ எனில், $x \in D$ அல்லது $x \in E$
  என்று காட்ட வேண்டும்.

  \begin{quote}
  இதுவரை வரையறைகளை மட்டுமே விரித்தோம்!{} $\subseteq$
  மற்றும் $\cup$ இன்றி முன்மொழிவை மீண்டும் கூறிவிட்டோம்;
  இப்போது அனைத்தையும் பற்றிய ஒரு நிபந்தனைக் கூற்றை
  நிறுவ வேண்டும். மேலே கண்டபடி, நிபந்தனையின் ``எனில்''
  பகுதி மெய்யான ஒரு $x$ ஐ எடுத்துக்கொண்டு, அதற்கு
  ``அப்போது'' பகுதியும் மெய் என்று காட்ட வேண்டும்.
  அதாவது $x \in B$ அல்லது $x \in C$ என்று
  எடுத்துக்கொண்டு, $x \in D$ அல்லது $x \in E$
  என்று காட்டுவோம்.\footnote{இந்தப் பத்தி நாம் செய்வதை
    விளக்குகிறது; இது நிறுவலின் பகுதி அல்ல. உங்கள் சொந்த
    நிறுவலை எழுதும்போது இவ்வளவு விவரிக்கத் தேவையில்லை.}
  \end{quote}

  $x \in B$ அல்லது $x \in C$ என்க. $x \in D$
  அல்லது $x \in E$ என்று காட்ட வேண்டும்.
  நிகழ்வுகளைப் பிரிக்கிறோம்.

  நிகழ்வு 1: $x \in B$. (c) படி $x \in D$. ஆகவே
  $x \in D$ அல்லது $x \in E$. (முந்தைய உட்பிரிவில்
  விவாதித்த உய்த்தறிதலையே இங்கு செய்தோம்!)

  நிகழ்வு 2: $x \in C$. (d) படி $x \in E$. ஆகவே
  $x \in D$ அல்லது $x \in E$.
\end{proof}

% SEGMENT OLP-0606-S12
\subsection{இருப்புக் கூற்றை நிறுவுதல்}

ஓர் இருப்புக் கூற்றை நிறுவச் சொன்னால், கேள்வி வழக்கமாக
``$\dots x \dots$ என்ற பண்புள்ள ஓர் $x$ உள்ளது என்று
நிறுவுக'' என்ற வடிவில் இருக்கும். அதாவது
``$\dots x \dots$'' என்று விவரிக்கப்பட்ட பண்புள்ள
ஒரு பொருள் இருக்கிறது என்று நிறுவ வேண்டும். அப்போது பொருத்தமான
பொருளைக் கண்டறிந்து, தேவையான பண்பு அதற்கு உள்ளது
என்று காட்ட வேண்டும். இது நேரடியாகத் தோன்றலாம்; ஆனால்
இத்தகைய நிறுவல் கடினமாக இருக்கலாம். ஒரு பொருளைக்
\emph{கட்டமைப்பது} அல்லது \emph{வரையறுப்பது}, பிறகு
அவ்வாறு வரையறுக்கப்பட்ட பொருளுக்கு வேண்டிய பண்பு
உள்ளது என்று நிறுவுவது வழக்கம். சரியான பொருளைக்
கண்டறிவதே கடினமாகலாம்; அதற்கு வேண்டிய பண்பு இருப்பதை
நிறுவுவதும் கடினமாகலாம்; சில நேரம் ஒரு பொருளை உண்மையில்
வரையறுத்துவிட்டோம் என்பதைக் காட்டுவதுகூட நுட்பமாக இருக்கலாம்!{}

% SEGMENT OLP-0606-S13
பொதுவாக, அந்தப் பொருளைக் குறிப்பிட்டு எழுதுவீர்கள்:
எடுத்துக்காட்டாக ``$x$ என்பது \dots'' என்று கூறலாம்.
இங்கு \dots{} நீங்கள் கருதும் பொருளைக் குறிப்பிடுகிறது.
தேவையானால் $\dots$ உண்மையில் இருக்கும் ஒரு பொருளை
விவரிக்கிறது என்று நிறுவி, பின்னர் $x$ க்கு $Q$
பண்பு உள்ளது என்று காட்டுவீர்கள். எளிய எடுத்துக்காட்டு:

\begin{prop}
  $x \in B$ என்க. அப்போது $A \subseteq B$ மற்றும்
  $A \neq \emptyset$ ஆகும் ஓர் $A$ உள்ளது.
\end{prop}

\begin{proof}
  $x \in B$ என்று எடுத்துக்கொள்க. $A = \{x\}$ என்க.
  \begin{quote}
    அதன் !!{element}s ஐப் பட்டியலிட்டு $A$ என்ற
    கணத்தை இங்கு வரையறுத்தோம். $x$ ஒரு பொருள் என்று
    எடுத்துக்கொண்டுள்ளோம்; அதன் !!{element}s ஐப்
    பட்டியலிட்டு ஒரு கணத்தை எப்போதும் அமைக்கலாம்.
    ஆகவே இங்கு ஒரு கணத்தை $A$ என்று வெற்றிகரமாக
    வரையறுத்தோம் என்பதைத் தனியாகக் காட்ட வேண்டியதில்லை.
    ஆனால் முன்மொழிவு கோரும் பண்புகள் $A$ க்கு உள்ளன
    என்பதைக் காட்ட வேண்டியுள்ளது; அதின்றி நிறுவல்
    நிறைவுபெறாது!{}
  \end{quote}
  $x \in A$ என்பதால் $A \neq \emptyset$.
  \begin{quote}
    இது $A$ ஐ $\{x\}$ என்று வரையறுத்ததையும்,
    $x \in \{x\}$ மற்றும் $x \notin \emptyset$
    என்ற வெளிப்படையான உண்மைகளையும் சார்ந்துள்ளது.
  \end{quote}
  $\{x\}$ இன் ஒரே !!{element} $x$; மேலும்
  $x \in B$. ஆகவே $A$ இன் ஒவ்வோர் !!{element}
  உம் $B$ இன் !!a{element}. $\subseteq$ இன்
  வரையறையின்படி $A \subseteq B$.
\end{proof}

% SEGMENT OLP-0606-S14
\subsection{இருப்புக் கூற்றுகளைப் பயன்படுத்துதல்}

ஓர் இருப்புக் கூற்று மெய் என்று உங்களுக்குத் தெரியும் என்க:
அதை நிறுவியிருக்கலாம் அல்லது பயன்படுத்தக்கூடிய எடுகோளாக
இருக்கலாம். எடுத்துக்காட்டாக, ``ஏதோ ஓர் $x$ க்கு
$x \in A$'' அல்லது ``ஓர் $x \in A$ உள்ளது''
என்க. அதை நிறுவலில் பயன்படுத்த, இருப்பதாகக் கூற்று
சொல்லும் பொருள்களில் ஒன்றுக்குப் பெயர் இருப்பதுபோல்
கருதலாம். $A$ இல் குறைந்தது ஒரு பொருள் இருப்பதால்,
அந்தப் பெயர் குறிக்கக்கூடிய பொருள் உள்ளது. அதை
குறிப்பாகத் தேர்ந்தெடுக்கவோ, $\in A$ என்பதற்கு மேல்
விவரிக்கவோ முடியாமல் இருக்கலாம். ஆயினும் நிறுவலுக்காக
அதைத் தேர்ந்தெடுத்ததுபோல் கருதிப் பெயரிடலாம்.
ஏற்கெனவே பயன்படுத்திய பெயரையோ எடுகோள்களில் வரும்
பெயரையோ தேர்ந்தெடுக்கக் கூடாது; இல்லையெனில்
தவறுகள் ஏற்படலாம். ``ஏதோ ஓர் $x$ க்கு $x \in A$''
என்பதிலிருந்து ``$a \in A$ என்க'' என்று மாறுவதன்
மூலம் இதைக் குறிக்கிறோம். இப்போது $a$ பற்றிச்
சிந்தித்து மற்ற எடுகோள்களைப் பயன்படுத்தலாம்; இறுதியில்
ஒரு முடிவு $p$ ஐ அடையலாம். அந்த $p$ இல் $a$
இனி வராவிட்டால், $p$ ஆனது $a \in A$ என்று
தற்காலிகமாக எடுத்துக்கொண்ட குறிப்பிட்ட பெயரைச் சாராது;
``ஏதோ ஓர் $x$ க்கு $x \in A$'' என்ற இருப்பு
எடுகோளிலிருந்தே அது பின்தொடர்கிறது.

% SEGMENT OLP-0606-S15
\begin{prop}
$A \neq \emptyset$ எனில் $A \cup B \neq \emptyset$.
\end{prop}

\begin{proof}
  $A \neq \emptyset$ என்க. ஆகவே ஏதோ ஓர் $x$
  க்கு $x \in A$.
  \begin{quote}
    இங்கு முதலில் முன்மொழிவின் எடுகோளை மீண்டும் கூறினோம்.
    $A \neq \emptyset$ என்ற எடுகோள் ஓர் இருப்புக்
    கூற்றை மறைத்திருக்கிறது; சில வரையறைகளை விரித்தால்தான்
    அது வெளிப்படும். $=$ இன் வரையறையின்படி,
    $A = \emptyset$ எனில், எனில் மட்டுமே, ஒவ்வொரு
    $x \in A$ உம் $\in \emptyset$ ஆகவும்,
    ஒவ்வொரு $x \in \emptyset$ உம் $\in A$ ஆகவும்
    இருக்கும். இருபுறத்தையும் மறுத்தால்,
    $A \neq \emptyset$ எனில், எனில் மட்டுமே,
    ஏதோ ஓர் $x \in A$ ஆனது $\notin \emptyset$
    அல்லது ஏதோ ஓர் $x \in \emptyset$ ஆனது
    $\notin A$. $\in \emptyset$ ஆக எதுவும்
    இல்லாததால் இரண்டாவது கூறு மெய்யாக முடியாது;
    ``$x \in A$ மற்றும் $x \notin \emptyset$''
    என்பதும் $x \in A$ என்பதாகவே சுருங்குகிறது.
    % SOURCE CORRECTION TA-MTH-002: இங்கு x அல்ல, கணம் A காலியற்றது.
    ஆகவே $A \neq \emptyset$ எனில், எனில் மட்டுமே,
    ஏதோ ஓர் $x$ க்கு $x \in A$. இது ஓர் இருப்புக்
    கூற்று. இப்போது $A$ இன் !!{element}s இல்
    ஒன்றுக்குப் பெயரிடுவதன் மூலம் அதைப் பயன்படுத்துகிறோம்:
  \end{quote}
  $a \in A$ என்க.
  \begin{quote}
    இப்போது $\in A$ ஆக உள்ள பொருள்களில் ஒன்றுக்குப்
    பெயர் அறிமுகப்படுத்தினோம். தொடர்ந்து $a$ பற்றிப்
    பேசுவோம்; ஆனால் $a \in A$ என்பதைத் தவிர வேறெதையும்
    எடுகோளாகக் கொள்ளாமல் கவனமாக இருப்போம்:
  \end{quote}
  $a \in A$ என்பதால் $\cup$ இன் வரையறையின்படி
  $a \in A \cup B$. ஆகவே ஏதோ ஓர் $x$ க்கு
  $x \in A \cup B$; அதாவது $A \cup B \neq \emptyset$.
  \begin{quote}
    கடைசிப் படியில் ``$a \in A \cup B$'' என்பதிலிருந்து
    ``ஏதோ ஓர் $x$ க்கு $x \in A \cup B$'' என்று
    மாறினோம். பிந்தைய கூற்றில் $a$ வரவில்லை; ஆகவே
    ``ஏதோ ஓர் $x$ க்கு $x \in A \cup B$'' என்பது
    ``ஏதோ ஓர் $x$ க்கு $x \in A$'' என்ற
    எடுகோளிலிருந்தே பின்தொடர்கிறது. அதுவே
    $A \cup B \neq \emptyset$ என்பதைக் குறிக்கிறது.
  \end{quote}
\end{proof}

% SEGMENT OLP-0606-S16
நாம் செய்ததுபோல் ``$x$'' என்ற கட்டுப்பட்ட மாறியையும்
$a$ போன்ற தற்காலிகப் பெயரையும் தனித்தனியாக வைத்திருப்பது
நல்ல பயிற்சி. ஆனால் நடைமுறையில் அவ்வாறு செய்யாமல்,
$x$ ஐயே பின்வருமாறு பயன்படுத்துவதுண்டு:
\begin{quote}
$A \neq \emptyset$ என்க; அதாவது ஓர் $x \in A$
உள்ளது. $\cup$ இன் வரையறையின்படி $x \in A \cup B$.
ஆகவே $A \cup B \neq \emptyset$.
\end{quote}
இவ்வாறு எழுதினால் வெவ்வேறு இருப்புக் கூற்றுகளுக்கு
வெவ்வேறு $x$, $y$ ஆகியவற்றைப் பயன்படுத்துகிறீர்களா என்பதில்
கூடுதல் கவனம் தேவை. எடுத்துக்காட்டாக,
``$A \neq \emptyset$ மற்றும் $B \neq \emptyset$
எனில் $A \cap B \neq \emptyset$'' என்ற கூற்றுக்கு
பின்வருவது \emph{சரியானதல்ல} (அந்தக் கூற்றே பொய்):
\begin{quote}
$A \neq \emptyset$ மற்றும் $B \neq \emptyset$ என்க.
ஆகவே ஏதோ ஓர் $x$ க்கு $x \in A$; மேலும் ஏதோ
ஓர் $x$ க்கு $x \in B$. $x \in A$ மற்றும்
$x \in B$ என்பதால், $\cap$ இன் வரையறையின்படி
$x \in A \cap B$. ஆகவே $A \cap B \neq \emptyset$.
\end{quote}
தவறான படி எங்கு நிகழ்கிறது என்று கண்டறிந்து, முடிவு
ஏன் பின்தொடரவில்லை என்பதை விளக்க முடியுமா?

\end{document}
